% Generated draft: requires mathematical and source review.
\documentclass[11pt]{article}
\usepackage[margin=25mm]{geometry}
\usepackage{fontspec}
\setmainfont{Latin Modern Roman}
\usepackage{amsmath,amssymb,mathtools,unicode-math}
\setmathfont{Latin Modern Math}
\usepackage{enumitem,needspace,xurl,hyperref}
\setlength{\parindent}{0pt}
\setlength{\parskip}{5pt}
\setlength{\emergencystretch}{3em}
\newcommand{\N}{\mathbb{N}}
\newcommand{\Z}{\mathbb{Z}}
\newcommand{\Q}{\mathbb{Q}}
\newcommand{\R}{\mathbb{R}}
\newcommand{\C}{\mathbb{C}}
\newcommand{\F}{\mathbb{F}}
\newcommand{\A}{\mathbb{A}}
\newcommand{\PP}{\mathbb{P}}
\newcommand{\E}{\mathbb{E}}
\newcommand{\eps}{\varepsilon}
\newcommand{\dd}{\,\mathrm{d}}
\begin{document}
\section*{UnsolvedMath ID: 11000193}
\textbf{Source problem code:} catalogue entry
\subsection{Definitions and mathematical statement}
Assume \(b=0\), so \(\Sigma\) is closed of genus \(g\). Let \(\pi,G\) and the component invariant \(e\) have their source meanings; none is defined in this extract. Conjecture 3.1 asks that for every integer
\[
 1\le k\le2g-2,
\]
the \(\operatorname{Mod}_\Sigma\)-action on
\[
 e^{-1}(2-2g+k)
\]
be ergodic. This simplifies the source's \(2-2g+b+k\) and \(2g+b-2\) expressions using its explicit \(b=0\) hypothesis, without proposing a version for nonzero \(b\).

For a specified invariant measure \(\mu\), ergodicity means every measurable invariant set has either measure zero or complement of measure zero. The record supplies no measure or measure class. It writes the components as subsets of \(\operatorname{Hom}(\pi,G)\), yet gives the identification
\[
 e^{-1}(3-2g)\approx\Sigma\times\R^{6g-8},
\]
a \((6g-6)\)-dimensional thickening. This is consistent with the dimension of a usual surface character-variety model, rather than the unquotiented representation space; the quotient and smooth/topological meaning of \(\approx\) require checking. A mapping-class action defined through outer automorphisms also needs this distinction. An Euler-class interpretation of \(e\), for example for \(\mathrm{PSL}_2(\R)\), is plausible terminology but is not a supplied choice of \(G\).

The context contrasts the action on this thickening with realization on \(\Sigma\) itself. It attributes to Morita [81] nonexistence of a smooth action inducing the natural projection \(\operatorname{Diff}(\Sigma)\to\operatorname{Mod}_\Sigma\), with precise regularity/genus hypotheses absent. It reports a then-recent announcement by Markovic [80] excluding analogous homeomorphism sections for closed genus greater than \(5\). These are source-era reports, not current verification. They do not supply the omitted measure or settle the ergodicity conjecture.
\subsection{Short English statement}
For a closed surface, the specified nonextremal components are conjectured ergodic under the mapping class action. The representation group, quotient and measure are missing. The source's thickening model and historical smooth/homeomorphism realization obstructions are retained as context.
\subsection{Sources}
\begin{itemize}
\item Local UnsolvedMath registry, permanent ID \(11000193\), exact source code \texttt{catalogue entry}, Conjecture 3.1, PDF page 220. No individual proposer or proposed year is supplied; the title is truncated.
\item Farb, edited volume \emph{Problems on Mapping Class Groups and Related Topics}: \url{https://www.math.uchicago.edu/~farb/papers/mcgbook.pdf}. The volume has 329 explicitly typeset numbered items across chapters, repeated local numbering distinguished by PDF page.
\item Internal references: Morita [81] for the reported smooth nonrealization and Markovic [80] for the source-era homeomorphism announcement in genus greater than five. Full entries, dates and exact theorem hypotheses are absent.
\item Extraction pdf-font-structure, accessed 2026-07-29. Source evidence NEEDS\_REVIEW presents an open problem and calls for a current primary-literature audit. Registry status open and classification OPEN-TRIAGE; the research summary says only ``Literature status: Truncated.''
\item Dataset attribution: Ulam AI / UnsolvedMath Contributors, CC BY 4.0. Source rights NEEDS\_REVIEW: publicly accessible author/publisher copy, redistribution terms requiring review.
\end{itemize}
\end{document}
